\documentclass[11pt,reqno]{amsart}
\usepackage[T1]{fontenc}\usepackage{lmodern}
\usepackage{amsmath,amssymb,amsthm}
\usepackage[margin=1in]{geometry}\usepackage{microtype}\usepackage{xcolor}
\usepackage[colorlinks=true,linkcolor=blue!55!black,citecolor=blue!55!black]{hyperref}
\usepackage{natbib}
\newcommand{\CM}{C_M}\newcommand{\phio}{\varphi^{\circ}}
\setlength{\parskip}{3pt}
\begin{document}
\title[Related work and positioning]{Related Work and Positioning}
\author{Team A1 --- Role 1}
\date{Supporting section}
\begin{abstract}
Where this paper sits, and---stated plainly---which of its results are already known. Two
of our five findings have close published antecedents; one is crowded; two appear to be
new. We say which is which here rather than let a referee discover it.
\end{abstract}
\maketitle

\section{Curtailable load contracts}
The instrument we study is a curtailment tariff, and the theory of selling curtailable
service is mature. \citet{ChaoWilson1987} established that a menu of contingent curtailment
contracts induces efficient self-selection and dominates random rationing; the market-structure
and investment extensions are \citet{ChaoOrenSmithWilson1988}. \citet{OrenSmith1992} is the
canonical operations treatment of designing and dispatching a curtailable-service program,
and is the closest existing analogue to our dispatch comparison. \citet{StraussOren1993} adds
notification lead time---relevant because we assume the load is instantaneously curtailable,
an assumption their results show is not innocuous. \citet{KamatOren2002} price interruptible
supply as a compound option and \citet{BaldickKolosTompaidis2006} derive a retailer's optimal
interruption policy under stochastic prices; neither considers emissions, which is the
opening we take. In our own venue, \citet{FattahiDasuAhmadi2023} is the current reference
point for direct load-control contract design, and is cost-focused with no carbon analysis.

Our surplus contract is a degenerate one-class priority contract, and our budget family
$d_\tau$ is a priority-service menu. We claim no novelty in the contract form. What we add
is the carbon ledger over that family and the instrument comparison of \S5.3.

\section{Flexible computing as a grid resource}
That computing load responds to electricity signals at all traces to \citet{Qureshi2009};
\citet{Liu2011} showed geographic load balancing can raise renewable utilization, though
under a fixed generation mix---the extension to endogenous renewable capacity is where our
capacity-substitution channel lives. \citet{Wierman2014} surveys operational feasibility, and
\citet{Radovanovic2023} documents a deployed carbon-shifting system at Google, which uses
marginal-intensity signals rather than a surplus rule. For mining specifically,
\citet{Menati2023aae} finds the carbon footprint of mining load varies up to 50\% from the
system average by location, and \citet{MenatiLeeXie2023} is the reference ERCOT simulation.
\citet{vanderJagt2024} models exactly our $K_M$ object as a ``demand sink'' and finds
low-capital-cost sinks cut curtailment substantially; \citet{RiepinBrown2024} raises the
additionality concern that bears on our Pareto framing.

\section{What is already known, stated plainly}
\textbf{Our negative result is not new.} \citet{KnittelSengaWang2025} find, empirically, that
flexible data centers flatten net demand, cut system costs, and \emph{raise} emissions when
the renewable share is low, with the sign turning favourable as penetration grows. That is
our ``the unconstrained rule raises carbon except in a cheap-renewable corner,'' in a
different currency. The general mechanism is older still: \citet{HollandMansur2008} showed
that reducing demand variance raises emissions in some regions and lowers it in others. We
therefore do not present the direction of the effect as a finding. What we contribute is the
\emph{exact analytic boundary} in primitives, where they give regional simulations and
reduced-form estimates.

This places our result in a recognized class. \citet{CarsonNovan2013} found bulk storage
arbitrage adds roughly $0.19$ tons CO\textsubscript{2} per MWh stored in ERCOT, and
\citet{HittingerAzevedo2015} confirmed the sign nationally while showing it flips as
renewable capacity rises---the empirical shadow of a classification condition.
\citet{Ding2023} shows curtailment caps can cut renewable capacity and raise
CO\textsubscript{2}, which is our feedback from surplus use to $K_R$ running in reverse.

\textbf{Our distribution-free statement is crowded, not pre-empted.} We have found no proof
of the Pareto claim for all $F$. But \citet{vanderJagt2024} simulates the mechanism,
\citet{Butters2025} gives the full structural equilibrium for a flexible asset in
\emph{Econometrica}, and \citet{HittingerAzevedo2015} documents the conditional flip. The
honest defense is not novelty of conclusion but of scope: distribution-freeness buys
robustness to exactly the heterogeneity that \citet{Novan2015},
\citet{CallawayFowlieMcCormick2018}, and \citet{Menati2023aae} document empirically---
identical resources, wildly different carbon signs by region and season.

\textbf{Our fossil-minimizing dispatch result is short and, we suspect, folk knowledge.} It
is a lemma, and we present it as one.

\section{Capacity investment}
Our three-capacity program is a peak-load-pricing and reliability-constrained capacity model
in the tradition of \citet{CrewKleindorfer1976}, surveyed for operations by
\citet{VanMieghem2003}. The distribution-free cutoffs are newsvendor critical fractiles:
$\phio=(\beta+p-c-C_F)/(\beta+p-c)$ \emph{is} a critical fractile, and we frame it as one.
\citet{Hu2015} warns that data granularity on intermittency drives renewable
over-investment---a caution about our own coarse residual-demand primitive.
\citet{AflakiNetessine2017} is the closest precedent for a perverse-sign result in this
literature: carbon pricing can discourage renewable investment under intermittency. That
paper is our evidence that ``green policy, wrong sign'' theorems publish in operations.
\citet{KokShangYucel2018} is the nearest \emph{Management Science} paper to our object---
pricing policy to renewable investment to emissions, with a sign that flips on market
characteristics---and \citet{KokShangYucel2020} establishes renewable/flexible
complementarity. \citet{PengWuSouza2024} asks our complementarity question for the
(renewable, flexible, storage) triple; since our third asset is a flexible \emph{load}
rather than storage, we owe a referee an explicit distinction, and the honest answer is
that storage is absent from our model and is the first extension a reader should demand.
\citet{SunarBirge2019} and \citet{Drake2016} set the modelling conventions we follow.

\section{Distribution-free methods}
Our ``for every $F$'' results sit in the tradition of \citet{Scarf1958} and
\citet{GallegoMoon1993}. \citet{PerakisRoels2008} formalizes what is and is not decidable
under partial distributional information, which is the right frame for our conditional-mean
dependence. We deliberately do \emph{not} claim a distributionally robust optimization
result: a referee would reasonably ask why we do not simply solve a DRO problem over $F$
\citep{DelageYe2010,WiesemannKuhnSim2014}, and our worst-case guarantee has infimum zero,
so the DRO vocabulary would promise more than we deliver. On detail-free design, our
contract is simple because the designer will not commit to a distribution, in the spirit of
\citet{Wilson1987} and \citet{BergemannMorris2005}; \citet{Carroll2015} is the model for how
to make that argument, and we note that our setting has \emph{no private information}, so we
avoid mechanism-design language entirely.

\section{Carbon accounting}
$\Phi$ is built on conditional means, which is marginal-emissions-factor logic done exactly;
\citet{SilerEvans2012} established that average factors badly misestimate avoided emissions.
Two boundaries of our accounting deserve statement. Our counterfactual is \emph{no load},
but a curtailed miner relocates rather than vanishing, so grid-local accounting understates
leakage \citep{Fowlie2009}; and our unconstrained-rule result is a supply-side backfire of
the kind catalogued by \citet{GillinghamRapsonWagner2016}. \citet{BushnellNovan2021} further
suggests our two-conditional-mean summary may be too coarse to capture ramping and start-up
effects.

\section{Summary of the contribution}
Against this literature we claim: (i) a tractable joint capacity characterization with
distribution-free cutoffs and a recursive-triangular uniqueness argument; (ii) an exact
closed-form carbon ledger on the large-\textsc{fcl} regime; and (iii) the price-versus-quantity
comparison of \S5.3---allocation equivalence at the greening thresholds, revenue divergence,
and a strict long-run ranking---which we have not found anywhere in this literature. We do
\emph{not} claim the direction of the naive-rule carbon effect, which is
\citet{KnittelSengaWang2025} and \citet{HollandMansur2008}.

\bibliographystyle{plainnat}
\bibliography{a1_references}
\end{document}
